% Part: sets-functions-relations
% Chapter: relations
% Section: operations

\documentclass[../../../include/open-logic-section]{subfiles}

\begin{document}

\olfileid{sfr}{rel}{ops}
\olsection{Operasi tumrap Relasi}

Asring migunani yen ngowahi utawa nggabungake relasi.
Ing \olref[sfr][rel][ord]{prop:stricttopartial},
kita ngrembug \emph{gabungan} relasi, kang mung
gabungan rong relasi nalika dianggep himpunan pasangan.
Semono uga, ing \olref[sfr][rel][ord]{prop:partialtostrict},
kita ngrembug selisih relatif relasi.
Ing ngisor iki ana sawetara operasi liyane
kang bisa ditindakake tumrap relasi.

\begin{defn}\ollabel{relationoperations}
Umpamakna $R$ lan $S$ iku relasi, lan $A$ iku himpunan apa wae.

\emph{Invers} saka $R$ yaiku
$R^{-1} = \Setabs{\tuple{y, x}}{\tuple{x,y} \in R}$.

\emph{Produk relatif} saka $R$ lan $S$ yaiku
$(R \mid S) =\{\tuple{x, z} : \exists y(Rxy \land Syz)\}$.

\emph{Watesan} saka $R$ ing $A$ yaiku
$\funrestrictionto{R}{A}= R\cap A^2$.

\emph{Aplikasi} saka $R$ marang $A$ yaiku
$\funimage{R}{A} = \{y :( \exists x \in A)Rxy\}$
\end{defn}

\begin{ex}
Umpamakna $S \subseteq \Int^2$ iku relasi panerus
ing~$\Int$, yaiku
$S = \Setabs{\tuple{x, y} \in \Int^2}{x + 1 = y}$,
mula $Sxy$ yen lan mung yen $x + 1 = y$.

$S^{-1}$ iku relasi pandhisik ing $\Int$, yaiku
$\Setabs{\tuple{x,y}\in\Int^2}{x -1 =y}$.

$S\mid S$ yaiku
$ \Setabs{\tuple{x,y}\in\Int^2}{x + 2 =y}$

$\funrestrictionto{S}{\Nat}$ iku relasi panerus ing~$\Nat$.

$\funimage{S}{\{1,2,3\}}$ yaiku $\{2, 3, 4\}$.
\end{ex}

\begin{defn}[Tutupan transitif]Umpamakna $R \subseteq A^2$ iku relasi biner.

\emph{Tutupan transitif} saka~$R$ yaiku
$R^+ = \bigcup_{0 < n \in\Nat} R^n$,
kanthi definisi rekursif $R^1 = R$ lan
$R^{n+1} = R^n\mid R$.

\emph{Tutupan refleksif transitif} saka $R$ yaiku
$R^* = R^+ \cup\Id{A}$.
\end{defn}

\begin{ex}
Jupuk relasi panerus $S \subseteq \Int^2$.
$S^2xy$ yen lan mung yen $x + 2 =y$,
$S^3xy$ yen lan mung yen $x + 3 = y$,
lan sapanunggalane.
Mula $S^+xy$ yen lan mung yen $x + n = y$
kanggo sawijining $n \geq 1$.
Kanthi tembung liya, $S^+xy$ yen lan mung yen
$x < y$, lan $S^*xy$ yen lan mung yen $x \le y$.
\end{ex}

\begin{prob}
Tuduhna yen tutupan transitif saka $R$ pancen transitif.
\end{prob}

\end{document}
